\documentclass[11pt]{article}
\usepackage{amsmath,amssymb,amsthm}
\usepackage[margin=1in]{geometry}
\usepackage{hyperref}
\usepackage{mdframed}

\newtheorem{theorem}{Theorem}[section]
\newtheorem{lemma}[theorem]{Lemma}
\newtheorem{corollary}[theorem]{Corollary}
\newtheorem{proposition}[theorem]{Proposition}
\theoremstyle{definition}
\newtheorem{definition}[theorem]{Definition}
\theoremstyle{remark}
\newtheorem{remark}[theorem]{Remark}

\newcommand{\R}{\mathbb{R}}
\newcommand{\norm}[1]{\lVert #1 \rVert}
\newcommand{\smin}{\sigma_{\min}}
\newcommand{\smax}{\sigma_{\max}}
\newcommand{\fnorm}[1]{\lVert #1 \rVert_F}
\newcommand{\tnorm}[1]{\lVert #1 \rVert_2}

\title{The Bischof--Stewart Pivoting Strategy:\\Mathematics, Complexity, and Pseudocode}
\author{Anil Damle}
\date{}

\begin{document}
\maketitle

\begin{abstract}
We give a self-contained exposition of the column pivoting strategy
based on incremental condition estimation, originally proposed by
Bischof~\cite{bischof1990incremental} and developed for column selection
by Stewart~\cite{stewart1998incremental}.  The key idea is to choose
pivots that greedily minimize the per-step growth of
$\fnorm{R^{-1}}$, where $R$ is the upper-triangular factor of the
evolving QR factorization.  We derive the incremental update formulas,
present detailed pseudocode using 1-indexed matrices, and analyze the
computational complexity.  Combined with Osinsky's
analysis~\cite{osinsky2025close}, the resulting algorithm selects $k$
columns from a $k \times n$ matrix with orthonormal rows in $O(nk^2)$
operations, achieving near-optimal submatrix conditioning guarantees.
\end{abstract}

%% ============================================================
\section{Introduction}
\label{sec:intro}
%% ============================================================

In the Golub--Klema--Stewart (GKS)
algorithm~\cite{golub_klema_stewart} for column subset selection, one
computes the leading $k$ right singular vectors $V_k \in \R^{n \times
k}$ of a matrix $A \in \R^{m \times n}$ and then selects $k$ columns
by applying a column-pivoted QR factorization to $V_k^T \in \R^{k
\times n}$.  Since $V_k$ has orthonormal columns, $V_k^T$ has
orthonormal rows, and column selection reduces to finding a $k \times
k$ submatrix $V_{J,1:k}$ with small $\fnorm{V_{J,1:k}^{-1}}$ (or
equivalently, large $\smin(V_{J,1:k})$).

The standard Businger--Golub~\cite{businger_linear_least_squares}
pivoting strategy selects the column of largest norm at each step.
While efficient, it does not directly control the conditioning of the
selected submatrix.  Bischof~\cite{bischof1990incremental} introduced
\emph{incremental condition estimation}, a technique that tracks the
growth of $\fnorm{R^{-1}}$ as columns are appended to the triangular
factor during QR factorization.
Stewart~\cite{stewart1998incremental} adapted this idea to column
selection: at each step, the pivot is chosen to minimize the
incremental growth of $\fnorm{R^{-1}}$.

Osinsky~\cite{osinsky2025close} proved that this greedy strategy,
when applied to a matrix with orthonormal rows, achieves
\begin{equation}\label{eq:osinsky-bounds}
  \fnorm{V_{J,1:k}^{-1}} \leq \sqrt{k(n-k+1)}
  \qquad\text{and}\qquad
  \tnorm{V_{J,1:k}^{-1}} \leq \sqrt{1 + k(n-k)},
\end{equation}
matching the asymptotic guarantees of Gu--Eisenstat strong rank-revealing
QR~\cite{gu_srrqr}.  These bounds feed directly into the GKS error
analysis: the spectral norm column selection error satisfies
$\tnorm{A - A_{:,J}(A_{:,J})^\dagger A} \leq \sigma_{k+1}(A)\sqrt{1 +
k(n-k)}$, while the Frobenius norm error satisfies $\fnorm{A -
A_{:,J}(A_{:,J})^\dagger A} \leq \fnorm{\Sigma_\bot}\sqrt{1 +
k(n-k)/r_k}$~\cite{armstrong2025structure}.

The purpose of this note is to carefully derive the mathematics
underlying the Bischof--Stewart criterion, present complete pseudocode
with 1-indexed matrices, and analyze the computational cost.

%% ============================================================
\section{Incremental Update of the Triangular Inverse}
\label{sec:incremental}
%% ============================================================

We work with a matrix $M \in \R^{k \times n}$ satisfying $MM^T =
I_k$ (in the GKS context, $M = V_k^T$).  A column-pivoted QR
factorization builds, over $k$ steps, a permutation $\Pi$ and
orthogonal $Q$ such that
\[
  QM\Pi = \begin{pmatrix} R \\ 0 \end{pmatrix}
\]
with $R \in \R^{k \times n}$ upper trapezoidal (here $R$ is $k \times
n$ since $M$ already has $k$ rows; the zero block is vacuous and $Q$
is $k \times k$ orthogonal).  More precisely, after $k$ Householder
steps one obtains $QM\Pi = R$ with $R$ upper trapezoidal.

We describe the state after $i - 1$ completed steps ($1 \leq i \leq
k$) and derive the formula governing the $i$-th pivot selection.

\subsection{State after \texorpdfstring{$i-1$}{i-1} steps}

After $i - 1$ steps, orthogonal transformations $Q_{i-1}$ have been
applied from the left and columns have been permuted so that
\begin{equation}\label{eq:partial-qr}
  Q_{i-1}^{\,T} M \Pi_{i-1}
  = \begin{pmatrix} R_{11} & R_{12} \\[2pt] 0 & \widetilde{M} \end{pmatrix},
\end{equation}
where $R_{11} \in \R^{(i-1) \times (i-1)}$ is upper triangular with
nonzero diagonal, $R_{12} \in \R^{(i-1) \times (n-i+1)}$, and
$\widetilde{M} \in \R^{(k-i+1) \times (n-i+1)}$ is the unreduced
block.  The first $i - 1$ columns (post-permutation) have been fully
reduced; the remaining $n - i + 1$ columns are candidates for the
$i$-th pivot.

\subsection{Extending the factorization by one column}

Suppose we select column $j$ from the unreduced block (i.e., the
$j$-th of the $n - i + 1$ remaining columns) as the next pivot.
Denote by
\[
  r_j = R_{12}(:,j) \in \R^{i-1}
  \qquad\text{and}\qquad
  \tilde{m}_j = \widetilde{M}(:,j) \in \R^{k-i+1}
\]
the top and bottom portions of this column
in~\eqref{eq:partial-qr}.  A Householder reflection $H_i$ is
constructed to zero out the entries of $\tilde{m}_j$ below its first
component, producing the new diagonal element
\begin{equation}\label{eq:rho}
  \rho_j = \tnorm{\tilde{m}_j}.
\end{equation}
After swapping column $j$ into position $i$ and applying $H_i$, the
triangular factor grows from $R_{11}$ to the $i \times i$ matrix
\begin{equation}\label{eq:R-extend}
  R_{11}^{(i)}
  = \begin{pmatrix} R_{11} & r_j \\[2pt] 0 & \rho_j \end{pmatrix}.
\end{equation}

\subsection{Incremental inverse formula}

The inverse of the extended triangular factor~\eqref{eq:R-extend} is
\begin{equation}\label{eq:R-inv-extend}
  \bigl(R_{11}^{(i)}\bigr)^{-1}
  = \begin{pmatrix} R_{11}^{-1} & -R_{11}^{-1} r_j / \rho_j \\[2pt]
                     0 & 1/\rho_j \end{pmatrix}
  = \begin{pmatrix} R_{11}^{-1} & -w_j / \rho_j \\[2pt]
                     0 & 1/\rho_j \end{pmatrix},
\end{equation}
where we define
\begin{equation}\label{eq:w-def}
  w_j = R_{11}^{-1} r_j \in \R^{i-1}.
\end{equation}
This is verified by direct multiplication:
$R_{11}^{(i)} \cdot (R_{11}^{(i)})^{-1}$ yields the identity.

\subsection{Growth of the inverse Frobenius norm}

Taking the squared Frobenius norm of~\eqref{eq:R-inv-extend}:
\begin{align}
  \fnorm{(R_{11}^{(i)})^{-1}}^2
  &= \fnorm{R_{11}^{-1}}^2
     + \frac{\tnorm{w_j}^2}{\rho_j^2}
     + \frac{1}{\rho_j^2} \notag \\
  &= \fnorm{R_{11}^{-1}}^2
     + \frac{1 + \tnorm{w_j}^2}{\rho_j^2}.
     \label{eq:frob-growth}
\end{align}
The second term is the \emph{increment} in $\fnorm{R^{-1}}^2$ caused
by appending column $j$.

\subsection{The Bischof--Stewart pivot criterion}

The Bischof--Stewart strategy~\cite{bischof1990incremental,
stewart1998incremental} selects the pivot that minimizes this
increment.  At step $i$, define the criterion for each remaining
column $j \in \{i, i+1, \ldots, n\}$ (after the permutation applied
so far):
\begin{equation}\label{eq:criterion}
  c_j
  = \frac{1 + \tnorm{w_j}^2}{\rho_j^2}
  = \frac{1 + \tnorm{R_{11}^{-1}\, R_{12}(:,j)}^2}
         {\tnorm{\widetilde{M}(:,j)}^2}.
\end{equation}
The pivot is
\begin{equation}\label{eq:pivot-rule}
  j^* = \operatorname*{argmin}_{j \in \{i,\ldots,n\}} c_j.
\end{equation}
This greedy rule ensures that each step contributes as little as
possible to $\fnorm{R^{-1}}^2$.

\begin{remark}[Interpretation]\label{rem:interp}
The numerator $1 + \tnorm{w_j}^2$ measures how much the existing
inverse ``resists'' the new column: if $r_j$ is nearly in the column
span of $R_{11}$, then $\tnorm{w_j}$ is large and the criterion
penalizes this choice.  The denominator $\rho_j^2 =
\tnorm{\tilde{m}_j}^2$ rewards columns whose unreduced portion is
large, as these produce a large diagonal element $\rho_j$ and hence a
small contribution $1/\rho_j^2$ to the inverse.  The criterion
balances these two effects.
\end{remark}

\begin{remark}[Relation to the spectral norm]\label{rem:spectral}
While the criterion directly controls $\fnorm{R^{-1}}$, the spectral
norm $\tnorm{R^{-1}}$ is also bounded.  Since $\tnorm{R^{-1}} \leq
\fnorm{R^{-1}}$, any bound on the Frobenius norm of the inverse
implies a (possibly looser) bound on the spectral norm.
Osinsky~\cite{osinsky2025close} establishes the sharper spectral bound
$\tnorm{R^{-1}} \leq \sqrt{1 + k(n-k)}$ directly.
\end{remark}

\subsection{Incremental maintenance of the \texorpdfstring{$w$}{w} vectors}
\label{sec:w-update}

A na\"ive implementation would recompute $w_j = R_{11}^{-1} r_j$ from
scratch at each step $i$ for every remaining column $j$, at a cost of
$O(i^2)$ per column via triangular solve.  Over all steps and columns,
this would total $O(nk^3)$---unacceptably expensive.  The key to an
efficient implementation is that the $w$ vectors can be updated
incrementally.

After step $i$ selects pivot $j^*$, the triangular factor grows
from $R_{11} \in \R^{(i-1)\times(i-1)}$ to $R_{11}^{(i)} \in
\R^{i \times i}$ as in~\eqref{eq:R-extend}.  For a remaining column
$\ell \neq j^*$, let $r_\ell^{(i)} \in \R^{i}$ denote its top $i$
entries in the updated partially reduced matrix (after the Householder
reflection $H_i$ has been applied).  Write
\[
  r_\ell^{(i)}
  = \begin{pmatrix} r_\ell \\ \alpha_\ell \end{pmatrix}
\]
where $r_\ell \in \R^{i-1}$ is the portion from the previous step and
$\alpha_\ell$ is the new $i$-th entry produced by $H_i$.  Then
\begin{equation}\label{eq:w-update}
  w_\ell^{(i)}
  = \bigl(R_{11}^{(i)}\bigr)^{-1} r_\ell^{(i)}
  = \begin{pmatrix} R_{11}^{-1} & -w_{j^*}/\rho_{j^*} \\
                     0 & 1/\rho_{j^*} \end{pmatrix}
    \begin{pmatrix} r_\ell \\ \alpha_\ell \end{pmatrix}
  = \begin{pmatrix} w_\ell - (\alpha_\ell / \rho_{j^*})\, w_{j^*} \\
                     \alpha_\ell / \rho_{j^*} \end{pmatrix},
\end{equation}
where $w_\ell = R_{11}^{-1} r_\ell$ is the $w$-vector from the
previous step.  The update~\eqref{eq:w-update} requires only:
\begin{enumerate}
\item the scalar $\alpha_\ell / \rho_{j^*}$ (one division, already
  available from the Householder step), and
\item a scaled vector subtraction $w_\ell \leftarrow w_\ell -
  (\alpha_\ell/\rho_{j^*})\,w_{j^*}$, costing $O(i)$ operations.
\end{enumerate}
Across all $n - i$ remaining columns, the update costs $O(i(n-i))$ at
step $i$, and summing over $i = 1, \ldots, k$ gives a total of
$O(nk^2)$ for maintaining all $w$-vectors.

Similarly, $\tnorm{w_\ell^{(i)}}^2$ can be maintained incrementally:
\begin{equation}\label{eq:wnorm-update}
  \tnorm{w_\ell^{(i)}}^2
  = \tnorm{w_\ell}^2
    - 2\frac{\alpha_\ell}{\rho_{j^*}}\, w_\ell^T w_{j^*}
    + \frac{\alpha_\ell^2}{\rho_{j^*}^2}\bigl(\tnorm{w_{j^*}}^2 + 1\bigr),
\end{equation}
which costs $O(i)$ per column (dominated by the inner product
$w_\ell^T w_{j^*}$, but this can be computed simultaneously with the
vector update).

%% ============================================================
\section{Algorithm}
\label{sec:algorithm}
%% ============================================================

Algorithm~\ref{alg:bs-qr} presents the complete Bischof--Stewart
column-pivoted QR factorization.  All indices are 1-based.  The
algorithm operates in-place on the input matrix $M$ and returns the
permutation vector $\pi$ identifying the selected columns.

\begin{figure}[ht]
\begin{mdframed}
\textbf{Algorithm 1:} Bischof--Stewart Column-Pivoted QR
\label{alg:bs-qr}

\medskip\noindent
\textbf{Input:} $M \in \R^{k \times n}$ with $MM^T = I_k$, \; $k \leq n$\\
\textbf{Output:} Permutation $\pi$ such that
  $\smin\!\bigl(M(:,\pi(1{:}k))\bigr)$ is controlled

\medskip\noindent
\begin{tabular}{@{}rp{0.88\textwidth}}
1: & $\pi \gets (1, 2, \ldots, n)$
     \hfill\textit{\% initialize permutation} \\
2: & $W \gets$ empty array of size $0 \times n$
     \hfill\textit{\% $W(:,j)$ stores $R_{11}^{-1}\, R_{12}(:,j)$} \\
3: & $s(j) \gets \tnorm{M(:,j)}^2$ for $j = 1, \ldots, n$
     \hfill\textit{\% squared column norms} \\[4pt]
4: & \textbf{for} $i = 1, \ldots, k$ \textbf{do} \\
   & \quad\textit{\% --- Pivot selection ---} \\
5: & \quad \textbf{for} $j = i, \ldots, n$ \textbf{do} \\
6: & \quad\quad \textbf{if} $i = 1$: \;
       $c(j) \gets 1 / s(j)$
       \hfill\textit{\% $w_j$ is empty} \\
7: & \quad\quad \textbf{else}: \;
       $c(j) \gets \bigl(1 + \tnorm{W(1{:}i{-}1,\,j)}^2\bigr) / s(j)$ \\
8: & \quad $j^* \gets \displaystyle\operatorname*{argmin}_{j \in \{i,\ldots,n\}} c(j)$ \\[4pt]
   & \quad\textit{\% --- Swap columns $i$ and $j^*$ ---} \\
9: & \quad Swap $\pi(i) \leftrightarrow \pi(j^*)$; \;
       swap $M(:,i) \leftrightarrow M(:,j^*)$ \\
10: & \quad Swap $W(:,i) \leftrightarrow W(:,j^*)$; \;
       swap $s(i) \leftrightarrow s(j^*)$ \\[4pt]
   & \quad\textit{\% --- Householder reduction of column $i$ ---} \\
11: & \quad Compute Householder $H_i$ s.t.\
       $H_i\, M(i{:}k,\,i) = \rho_i\, e_1$, \;
       $\rho_i = \tnorm{M(i{:}k,\,i)}$ \\
12: & \quad $M(i{:}k,\, i{:}n) \gets H_i\, M(i{:}k,\, i{:}n)$ \\[4pt]
   & \quad\textit{\% --- Update $W$ and $s$ for remaining columns ---} \\
13: & \quad \textbf{if} $i < k$ \textbf{then} \\
14: & \quad\quad $w^* \gets W(1{:}i{-}1,\, i)$
       \hfill\textit{\% $w$-vector of chosen pivot} \\
15: & \quad\quad \textbf{for} $j = i+1, \ldots, n$ \textbf{do} \\
16: & \quad\quad\quad $\alpha_j \gets M(i,j)$; \;\;
       $\beta_j \gets \alpha_j / \rho_i$ \\
17: & \quad\quad\quad \textbf{if} $i = 1$: \;
       $W(1,j) \gets \beta_j$ \\
18: & \quad\quad\quad \textbf{else}: \;
       $W(1{:}i{-}1,\, j) \gets W(1{:}i{-}1,\, j) - \beta_j \cdot w^*$; \;\;
       $W(i,j) \gets \beta_j$ \\
19: & \quad\quad\quad $s(j) \gets s(j) - \alpha_j^2$
       \hfill\textit{\% downdate column norm} \\
20: & \textbf{return} $\pi$
\end{tabular}
\end{mdframed}
\end{figure}

\begin{remark}[Column norm downdating]\label{rem:downdate}
Line~19 uses the standard norm downdating identity
$\tnorm{M(i{+}1:k,j)}^2 = \tnorm{M(i:k,j)}^2 - M(i,j)^2$.  In
floating-point arithmetic this subtraction can suffer cancellation
when $M(i,j)^2 \approx \tnorm{M(i:k,j)}^2$.  The classical remedy,
as in the Businger--Golub
algorithm~\cite{businger_linear_least_squares}, is to recompute the
column norm from scratch whenever the downdated value drops below a
threshold (e.g., when more than half the original norm has been
subtracted).  This safeguard adds at most $O(nk)$ operations over the
entire algorithm and does not change the asymptotic cost.
\end{remark}

\begin{remark}[First step simplification]\label{rem:first-step}
At $i = 1$, the triangular factor $R_{11}$ is empty and $w_j$ is
vacuous, so the criterion reduces to $c(j) = 1/s(j)$.  Minimizing
$c(j)$ is equivalent to maximizing $s(j) = \tnorm{M(:,j)}^2$, which
coincides with the Businger--Golub rule.  Thus the first pivot of the
Bischof--Stewart strategy agrees with standard column-pivoted QR\@.
\end{remark}

\begin{remark}[Tracking the full inverse norm]\label{rem:track-frob}
If the running value of $\fnorm{R^{-1}}^2$ is desired (e.g., for
diagnostic purposes), it can be maintained by initializing $F \gets 0$
and updating $F \gets F + c(j^*)$ after each pivot selection.  After
all $k$ steps, $F = \fnorm{R^{-1}}^2$ where $R = R_{11}^{(k)}$ is
the final $k \times k$ triangular factor.
\end{remark}

%% ============================================================
\section{Computational Complexity}
\label{sec:complexity}
%% ============================================================

We analyze the cost of Algorithm~\ref{alg:bs-qr} step by step.
Throughout, we assume $k \leq n$ and focus on the regime $k \ll n$
relevant to the GKS application.

\subsection{Per-step costs}

At step $i$ (selecting the $i$-th pivot), the dominant operations are:

\begin{enumerate}
\item \textbf{Criterion evaluation} (lines 6--10):
  Computing $c(j)$ for each of the $n - i + 1$ remaining columns
  requires reading $\tnorm{W(1:i-1,j)}^2$ and $s(j)$.  If
  $\tnorm{W(1:i-1,j)}^2$ is maintained as a running scalar (updated
  alongside $W$ via~\eqref{eq:wnorm-update}), each evaluation costs
  $O(1)$, for a total of $O(n - i)$ at step $i$.

\item \textbf{Column swaps} (lines 13--16):
  Swapping columns in $M$, $W$, $s$, and $\pi$ costs $O(k)$ for $M$,
  $O(i)$ for $W$, and $O(1)$ for $s$ and $\pi$; total $O(k)$.

\item \textbf{Householder reflection} (lines 18--19):
  Constructing $H_i$ costs $O(k - i)$.  Applying $H_i$ to the
  submatrix $M(i:k,\; i:n)$ of size $(k - i + 1) \times (n - i + 1)$
  costs $O((k - i)(n - i))$.

\item \textbf{$W$-vector update} (lines 22--30):
  For each of the $n - i$ remaining columns, the update~\eqref{eq:w-update}
  costs $O(i)$ (a scaled vector subtraction of length $i-1$ plus one
  scalar division).  Total: $O(i(n - i))$.

\item \textbf{Norm downdating} (line 31):
  One subtraction per remaining column: $O(n - i)$.
\end{enumerate}

\subsection{Total cost}

Summing over $i = 1, \ldots, k$:

\paragraph{Householder reflections.}
\[
  \sum_{i=1}^{k} O\bigl((k-i)(n-i)\bigr)
  = O\!\left(n\sum_{i=1}^{k}(k-i)\right)
  = O\!\left(n \cdot \frac{k(k-1)}{2}\right)
  = O(nk^2).
\]

\paragraph{$W$-vector maintenance.}
\[
  \sum_{i=1}^{k} O\bigl(i(n-i)\bigr)
  = O\!\left(n\sum_{i=1}^{k} i - \sum_{i=1}^{k} i^2\right)
  = O\!\left(n \cdot \frac{k(k+1)}{2}\right)
  = O(nk^2).
\]

\paragraph{Criterion evaluation and norm downdating.}
\[
  \sum_{i=1}^{k} O(n - i) = O(nk).
\]

\paragraph{Column swaps.}
\[
  \sum_{i=1}^{k} O(k) = O(k^2).
\]

The total cost is therefore
\begin{equation}\label{eq:total-cost}
  \boxed{T(n, k) = O(nk^2)}
\end{equation}
operations, dominated equally by the Householder reflections and the
$W$-vector updates.

\subsection{Comparison with alternative strategies}

\begin{enumerate}
\item \textbf{Businger--Golub column-pivoted
  QR}~\cite{businger_linear_least_squares}: Also $O(nk^2)$, but uses
  maximum-column-norm pivoting, which does not directly control
  $\fnorm{R^{-1}}$ and hence provides weaker conditioning guarantees.

\item \textbf{Gu--Eisenstat strong RRQR}~\cite{gu_srrqr}: Achieves
  optimal bounds on the entries of $R_{11}^{-1}R_{12}$, but the
  algorithm involves a post-processing phase of column swaps.  Each
  swap requires re-triangularizing via Givens rotations, and the
  number of swaps is bounded but can be substantial.  The worst-case
  cost exceeds $O(nk^2)$ when many swaps are needed.  In contrast,
  Bischof--Stewart achieves comparable asymptotic guarantees in a
  single forward pass.

\item \textbf{Standard partial QR (no pivoting)}: $O(nk^2)$, but
  without any pivoting the selected submatrix may be arbitrarily
  ill-conditioned.
\end{enumerate}

\begin{remark}[Practical considerations]\label{rem:practical}
The $W$-vector updates involve BLAS-1 operations (scaled vector
addition) that are memory-bandwidth limited on modern hardware.  In
practice, the updates can be organized as a matrix operation: defining
$W_{\mathrm{rem}} \in \R^{(i-1) \times (n-i)}$ as the block of
$w$-vectors for remaining columns, the update at step $i$ is a rank-1
modification $W_{\mathrm{rem}} \leftarrow W_{\mathrm{rem}} -
(\beta/\rho_i)\,w^* \cdot \alpha^T$, where $\alpha \in \R^{n-i}$
collects the new row entries.  This enables the use of BLAS-2
routines, improving cache utilization.
\end{remark}

%% ============================================================
\section{Connection to GKS Error Bounds}
\label{sec:connection}
%% ============================================================

The output of Algorithm~\ref{alg:bs-qr} is a permutation $\pi$ whose
first $k$ entries identify the selected column indices $J = \{\pi(1),
\ldots, \pi(k)\}$.  By Osinsky's
analysis~\cite{osinsky2025close}, the selected submatrix
$V_{J,1:k}$ satisfies the bounds~\eqref{eq:osinsky-bounds}.

These bounds enter the GKS error analysis through the master bound: for
any orthonormal $Z \in \R^{n \times k}$ approximating the leading
right singular subspace of $A$ and any index set $J$ with $|J| = k$,
\[
  \tnorm{A - A_{:,J}(A_{:,J})^\dagger A}
  \;\leq\;
  \frac{\tnorm{A - AZZ^T}}{\smin(Z_J)}.
\]
Setting $Z = V_k$ (exact singular vectors) and applying the
Bischof--Stewart selection yields
$\smin(V_{J,1:k}) \geq 1/\sqrt{1 + k(n-k)}$, giving the spectral
bound $\tnorm{A - A_{:,J}(A_{:,J})^\dagger A} \leq
\sigma_{k+1}(A)\sqrt{1 + k(n-k)}$.  The Frobenius bound follows
similarly from
$\fnorm{V_{J,1:k}^{-1}}^2 \leq k(n-k+1)$~\cite{armstrong2025structure}.

For the randomized GKS algorithm, the same column selection procedure
is applied to approximate right singular vectors $\hat{V}_k$ obtained
via randomized SVD\@.  Since $\hat{V}_k$ also has orthonormal columns,
Algorithm~\ref{alg:bs-qr} applies with identical guarantees on
$\smin(\hat{V}_{k,J})$, and the overall error bound involves an
additional factor from the subspace approximation
quality~\cite{halko_finding_structure_with_randomness}.

\bibliographystyle{plain}
\bibliography{gks}

\end{document}
